\begin{pastexam}[label={px:ncku-113-prog-12}]{成大資訊聯招 113 學年度 第 12 題}
（計算）

(10\%) Consider a neural network with $n$ layers of fully connected layers, where the first layer
is the input layer, the $n$th layer is the output layer, and the rest are hidden layers. The $i$th
layer has $P_i$ neurons, where $1\le i\le n$. Each neuron in the $i$th layer is connected to every
neuron in the $(i+1)$th layer, where $1\le i<n$. Let $w^i_{j,k}$ represent the weight of the
connection from the $j$th neuron in the $i$th layer to the $k$th neuron in the $(i+1)$th layer.
Let $[v^i_1\ v^i_2\ \cdots\ v^i_{P_i}]^T$ be the values of neurons in the $i$th layer, where
$1\le i\le n$. Then, we can obtain $v^{i+1}_k=\sum_{j=1}^{P_i} w^i_{j,k}v^i_j$, where
$1\le k\le P_{i+1}$. Given that $\langle P_1,P_2,P_3,P_4,P_5,P_6\rangle=\langle
10,6,12,5,50,3\rangle$, we want to achieve the result from the input layer to the output layer in
the fastest possible way. Please determine the minimum number of multiplications required for this
neural network.
\end{pastexam}

\begin{pastsoln}
答案：$1326$ 次乘法，順序是 $W^1\bigl(W^2\bigl(W^3(W^4 W^5)\bigr)\bigr)$。

\textbf{解讀}。 把第 $i$ 層到第 $i + 1$ 層的權重排成 $P_i \times P_{i+1}$ 的矩陣
$W^i = [w^i_{j,k}]$，題目的式子就是 $(v^{i+1})^T = (v^i)^T W^i$，所以
\[
(v^6)^T = (v^1)^T\, W^1 W^2 W^3 W^4 W^5 ,
\]
五個矩陣的維度依序是 $10 \times 6$、$6 \times 12$、$12 \times 5$、$5 \times 50$、$50 \times 3$。
題目只給層的大小、問「最快」的做法，考的是矩陣鏈乘（資結書第~18.2~節）：選一個乘法順序，讓算出
$W = W^1 W^2 W^3 W^4 W^5$（$10 \times 3$）的純量乘法次數最少；有了 $W$，輸入層到輸出層只要一次
$10 \times 3$ 的乘法。 下面照這個解讀計算。 （另一種讀法是只算\textbf{一筆}輸入：把 $(v^1)^T$ 當成
$1 \times 10$ 的矩陣放在鏈的最前面，最好的順序是由左往右一層一層算，共
$10 \cdot 6 + 6 \cdot 12 + 12 \cdot 5 + 5 \cdot 50 + 50 \cdot 3 = 592$ 次。 作答時寫出自己用哪一種讀法。）

\textbf{遞迴式}（資結書虛擬碼~18.4）。 維度 $p_0, \ldots, p_5 = 10, 6, 12, 5, 50, 3$，$W^i$ 是
$p_{i-1} \times p_i$。 $m[i, j]$ 是算出 $W^i \cdots W^j$ 最少要幾次乘法：
\[
m[i, i] = 0, \qquad
m[i, j] = \min_{i \le k < j} \bigl( m[i, k] + m[k + 1, j] + p_{i-1} p_k p_j \bigr).
\]

\textbf{長度 2}：
$m[1,2] = 10 \cdot 6 \cdot 12 = 720$，
$m[2,3] = 6 \cdot 12 \cdot 5 = 360$，
$m[3,4] = 12 \cdot 5 \cdot 50 = 3000$，
$m[4,5] = 5 \cdot 50 \cdot 3 = 750$。

\textbf{長度 3}：
\begin{align*}
m[1,3] &= \min(m[1,2] + 10 \cdot 12 \cdot 5,\ m[2,3] + 10 \cdot 6 \cdot 5) = \min(1320, 660) = 660
  \quad (k = 1), \\
m[2,4] &= \min(m[3,4] + 6 \cdot 12 \cdot 50,\ m[2,3] + 6 \cdot 5 \cdot 50) = \min(6600, 1860) = 1860
  \quad (k = 3), \\
m[3,5] &= \min(m[4,5] + 12 \cdot 5 \cdot 3,\ m[3,4] + 12 \cdot 50 \cdot 3) = \min(930, 4800) = 930
  \quad (k = 3).
\end{align*}

\textbf{長度 4}：
\begin{align*}
m[1,4] &= \min(m[2,4] + 3000,\ m[1,2] + m[3,4] + 6000,\ m[1,3] + 2500) \\
  &= \min(4860,\ 9720,\ 3160) = 3160 \quad (k = 3), \\
m[2,5] &= \min(m[3,5] + 216,\ m[2,3] + m[4,5] + 90,\ m[2,4] + 900) \\
  &= \min(1146,\ 1200,\ 2760) = 1146 \quad (k = 2).
\end{align*}
（$k = 1, 2, 3$ 的 $p_0 p_k p_4$ 是 $10 \cdot 6 \cdot 50$、$10 \cdot 12 \cdot 50$、$10 \cdot 5 \cdot 50$；
$k = 2, 3, 4$ 的 $p_1 p_k p_5$ 是 $6 \cdot 12 \cdot 3$、$6 \cdot 5 \cdot 3$、$6 \cdot 50 \cdot 3$。）

\textbf{長度 5}：
\begin{align*}
m[1,5] &= \min\bigl(m[2,5] + 10 \cdot 6 \cdot 3,\ m[1,2] + m[3,5] + 10 \cdot 12 \cdot 3, \\
  &\qquad\quad m[1,3] + m[4,5] + 10 \cdot 5 \cdot 3,\ m[1,4] + 10 \cdot 50 \cdot 3\bigr) \\
  &= \min(1146 + 180,\ 720 + 930 + 360,\ 660 + 750 + 150,\ 3160 + 1500) \\
  &= \min(1326,\ 2010,\ 1560,\ 4660) = 1326 \quad (k = 1).
\end{align*}

整理成表（上三角，格子裡是 $m[i, j]$，括號是最好的 $k$）：
\begin{center}
\small
\begin{tabular}{c|ccccc}
$i \backslash j$ & 1 & 2 & 3 & 4 & 5 \\
\hline
1 & 0 & 720 (1) & 660 (1) & 3160 (3) & 1326 (1) \\
2 &   & 0 & 360 (2) & 1860 (3) & 1146 (2) \\
3 &   &   & 0 & 3000 (3) & 930 (3) \\
4 &   &   &   & 0 & 750 (4) \\
5 &   &   &   &   & 0 \\
\end{tabular}
\end{center}
從 $m[1,5]$ 的 $k = 1$ 往回找：$W^1 (W^2 \cdots W^5)$；$m[2,5]$ 的 $k = 2$：$W^2 (W^3 W^4 W^5)$；
$m[3,5]$ 的 $k = 3$：$W^3 (W^4 W^5)$。 所以順序是 $W^1(W^2(W^3(W^4 W^5)))$，最少 $1326$ 次乘法。
直覺上，從最右邊開始乘，結果一直是「幾列 $\times$ 3 行」的窄矩陣，每一步都便宜；若先算
$W^3 W^4$，會產生 $12 \times 50$ 的大矩陣，光這一步就要 $3000$ 次。
\end{pastsoln}
